\documentclass[exercice,a4]{thedocument}%
\usepackage{texmaths-tikz}%
\usepackage{pstricks}%
\usepackage{pst-bezier}%
\usepackage{pstricks-add}%
\usepackage{pst-eucl}%
\usepackage{pst-func}%
\usepackage{pst-tree}%
\usepackage{pst-circ}%
\usepackage{tikz}%
\usepackage{pgfplots}%
\usepackage{tkz-tab}%
\usepackage{tikzpagenodes}%
\usepackage{tabularray}%
\usepackage{tabularx}%
\usepackage{multirow}%
\usepackage{multicol}%
\usepackage{scratch3}%
\usepackage{enumitem}%
\usepackage[np]{numprint}%
\usepackage{pas-tableur}%
\usepackage{listings}%
\usepackage{graphicx}%
\usepackage{adjustbox}%
\usepackage{needspace}%
\usepackage{xcolor}%
\usepackage{eurosym}%
\usepackage{diagbox}%
\usepackage{colortbl}%
\usepackage{multido}%
\usepackage{icomma}%
\usepackage[french]{babel}%
\graphicspath{{../brevet/2026/}}%
\startdatakeys
    \source{26GENMATAN1}
    \BO{2026}
    \cycle{4}
    \competencesDuSocle{CA,RA,CO}
    \objectifsApprentissage{C4A1aa,C4A2ad,C4B2cb,C4D2aa,C4D3ac}
    \motsCles{pourcentage,probabilité,conversion,puissance}
\enddatakeys
\definecolor{gris}{gray}{0.5}%
\providecommand{\wideparen}[1]{\overset{\frown}{#1}}%
\begin{document}%
\providecommand{\degres}{^\circ}%
Dans cet exercice, les deux parties sont indépendantes.

Une entreprise développe une intelligence artificielle (IA) capable de reconnaître des objets sur des images.

\subsubsection*{Partie A}

On entraîne l'IA à partir d'une base de données de \np{50000} images réparties en 4 catégories : \og Objets du quotidien \fg, \og Animaux \fg, \og Véhicules \fg, \og Autres \fg.

L'intelligence artificielle est testée pour mesurer sa précision et son efficacité. Les images sont réparties comme suit :

\vspace{6pt}
\begin{center}
\begin{tabular}{|>{\columncolor{gris}}c|l|c|}\hline
\rowcolor{gris}
& \textbf{Type d'image} & \textbf{Nombre d'images} \\\hline
2 & Objets du quotidien & 28\,000 \\\hline
3 & Animaux & \np{12000} \\\hline
4 & Véhicules & \np{8000} \\\hline
5 & Autres & ? \\\hline
\end{tabular}
\end{center}

\medskip

\begin{enumerate}
\item Combien d'images appartiennent à la catégorie \og Autres \fg\,?
\item Sur l'ensemble des tests, l'intelligence artificielle reconnaît correctement 90\,\% des \og Objets du quotidien \fg.

Calculer le nombre d'images reconnues correctement dans cette catégorie.
\item L'intelligence artificielle reconnaît correctement \np{5600} images de la catégorie \og Véhicules \fg.

Quel pourcentage de réussite cela représente-t-il dans cette catégorie\,?
\item Une image est tirée au hasard dans la base de données.\\
Quelle est la probabilité que l'image tirée soit l'image d'un \og Objet du quotidien \fg\,?

On donnera le résultat sous la forme d'un nombre décimal.
\end{enumerate}

\bigskip

\subsubsection*{Partie B}

L'intelligence artificielle, très utilisée dans le monde entier, nécessite une quantité importante d'électricité. L'énergie consommée peut s'exprimer en wattheures (Wh).

En 2024, sa consommation annuelle est estimée à 82\,000 Gigawattheures (GWh).

En comparaison, un collège consomme en moyenne 200\,000 kilowattheures (kWh) par an.

\medskip

\begin{center}
\begin{tabularx}{0.35\linewidth}{|X|}\hline
\textbf{Rappels :}\\
-- 1\,kWh $= 10^3$\,Wh\\
-- 1\,GWh $= 10^9$\,Wh\\ \hline
\end{tabularx}
\end{center}

\medskip

\begin{enumerate}[resume]
\item Convertir la consommation de l'IA et d'un collège en Wh.

\emph{Exprimer ces résultats sous la forme d'une écriture scientifique.}
\item Combien de collèges pourrait-on alimenter pendant un an avec la consommation électrique de l'intelligence artificielle\,?
\item En France, il y a environ \np{7100} collèges. Dans cette question, on suppose que chaque collège a la même consommation d'énergie annuelle moyenne (\np{200000} kWh).

Pendant combien d'années environ pourrait-on alimenter tous les collèges français avec la consommation électrique annuelle de cette intelligence artificielle\,?
\end{enumerate}

%\newpage

%%%%%%%%%%%%%%
\medskip
\showthesource
\end{document}
